% Chapter 1: Chapter 1. Introduction


\section{Background}\label{sec:1_1}

Convolutional Neural Networks (CNNs) have become the dominant approach for computer vision tasks including image classification, object detection, and semantic segmentation~\cite{SZE_SURVEY}. However, the computational intensity of CNN inference poses significant challenges for deployment on resource-constrained edge devices. A single forward pass through even a modest network like LeNet-5 requires millions of multiply-accumulate (MAC) operations, making general-purpose CPU execution prohibitively slow and power-inefficient for real-time applications~\cite{KUNG_SYSTOLIC}.


Hardware acceleration has emerged as the primary solution to this challenge. Application-Specific Integrated Circuits (ASICs) and Field-Programmable Gate Arrays (FPGAs) can achieve orders of magnitude better performance and energy efficiency than CPUs by exploiting the inherent parallelism in CNN computations. Systolic array architectures, in particular, have proven highly effective for matrix multiplication—the core operation in both fully-connected and convolutional layers—by organizing processing elements in a regular grid with local data communication.


\section{RISC-V and Custom Instruction Extensions}\label{sec:1_2}

RISC-V is an open standard instruction set architecture (ISA) that has gained significant traction in both academic and industrial settings since its introduction at UC Berkeley in 2010~\cite{RISCV_UNPRIV_SPEC,WATERMAN_PHD,RISCV_READER}. Unlike proprietary ISAs, RISC-V’s modular design explicitly reserves encoding space for custom instruction extensions, enabling domain-specific acceleration without breaking compatibility with the standard ecosystem.


The Nuclei Instruction Co-extension (NICE) interface, implemented in the E203 Hummingbird v2 processor~\cite{NUCLEI_E203}, provides a standardized mechanism for integrating custom hardware accelerators directly into the processor pipeline~\cite{NUCLEI_NICE}. Through NICE, custom instructions appear as native RISC-V instructions to the programmer, with the processor automatically handling operand fetch, hazard detection, and result writeback. This approach eliminates the overhead of memory-mapped I/O or coprocessor communication protocols, enabling fine-grained, low-latency acceleration.


\section{Problem Statement}\label{sec:1_3}

While RISC-V custom instruction extensions offer a promising path for CNN acceleration, the practical challenges of integrating a hardware accelerator with a RISC-V core and validating the complete system on an FPGA platform remain significant. Key challenges include:


Correctly implementing the NICE interface protocol for custom instruction handshaking and data transfer.


Managing the FPGA build flow, including block RAM initialization, clock domain crossing, and timing closure.


Establishing a reliable debugging methodology using Integrated Logic Analyzer (ILA) probes for observing internal SoC signals.


Bridging the gap between RTL simulation and hardware behavior, where synthesis-specific issues, memory initialization formats, and macro visibility can cause divergent behavior.


\section{Project Objectives}\label{sec:1_4}

This project aims to design, implement, and validate a lightweight CNN accelerator integrated with a RISC-V E203 processor through the NICE interface, targeting FPGA prototype validation on the Davinci Pro A7-100T development board.


The specific objectives are:


Design a 4×4 systolic PE array supporting INT8 quantized convolution, with six custom NICE instructions (CFG, CLEAR, WLOAD, DLOAD, COMP, RSTAT) for software control.


Integrate the accelerator with the E203 Hummingbird v2 SoC, including ITCM/DTCM memory, UART, and GPIO peripherals.


Establish a complete FPGA bring-up pipeline from RTL simulation through Vivado synthesis, place-and-route, and bitstream generation.


Validate the CPU boot chain by running a minimal bare-metal program (hello\_e203) on the FPGA board with UART output verification.


Test the CNN accelerator custom instructions on the FPGA using custom NICE test programs, with ILA-based diagnostic evidence collection.


\section{Thesis Structure}\label{sec:1_5}

Chapter~2 establishes the technical context, covering the RISC-V ISA, the E203 processor microarchitecture, CNN fundamentals, and FPGA prototyping methods. Chapter~3 details the system design: the SoC architecture, the NICE custom instruction set, the PE array implementation, and the FPGA bring-up strategy. Chapter~4 reports the experimental outcomes across RTL simulation, bitstream construction, board-level validation of the hello\_e203 program, root-cause analysis of CPU boot failures, and NICE accelerator testing with ILA evidence. Chapter~5 interprets these results, compares achievements against the project objectives, and identifies limitations. Chapter~6 summarizes the contributions and outlines directions for future work.


